\documentclass[10pt]{beamer}

\usepackage{verbatim}
%\usepackage[latin1]{inputenc}
\usepackage{amssymb,amsmath,amsfonts,url}
\usepackage{graphicx,color}
\usepackage{pgf,pgfarrows,pgfnodes,pgfautomata,pgfheaps,pgfshade}
\usepackage[utf8]{inputenc}
\usepackage[french]{babel}
\usepackage{tikz}

\usepackage[all]{xy}
\usepackage{amstext}
\usepackage{xspace}
\usepackage{mdwtab}


\mode<article> {
  \usepackage{fullpage}
  \usepackage{hyperref}
  \setjobnamebeamerversion{advancedcourse.beamer}
}





\newcommand{\alertonce}[1]{\addtocounter{beamerpauses}{-1}\alert<+>{#1}}
\newcommand{\alerttwice}[1]{\addtocounter{beamerpauses}{-1}\alert<+-+(1)>{#1}}

%\newcommand{\subsubsubsection}[1]{\noindent\textbf{#1}}

\newcommand{\cfbox}[1]{\begin{center}\fbox{#1}\end{center}}
\newcommand{\cbox}[1]{\begin{center} #1 \end{center}}

%%%%%%%%%%%%%%%%%%%%%%%%%%%%%%%%%%%%%%%%%%%%%%%%%%%%%%%%%%%%%%%%%%%%%%%%%%%%%%%%%%%%%%%%%

\graphicspath{{Graphics/}}
%\input def_0.tex

%\input{StyleSlides.tex}
\mode<presentation> {
    \usepackage{beamerthemeBoadilla}

  %\beamertemplatetransparentcovereddynamic
  \beamertemplateballitem
}


\title[IN100 Cours 6]{\textbf{Fondements informatiques I}\\
Cours 6: Listes}

\author[]{Sorina Ionica \texttt{sorina.ionica@uvsq.fr} \\~\\Sandrine Vial  \texttt{sandrine.vial@uvsq.fr} }


\institute{}

\date{}


%\setbeamercolor{normal text}{fg=black}

\begin{document}


\frame{\titlepage}


\begin{frame}[fragile]{Les listes}

Un \textcolor{blue}{tableau} est une structure de données contenant une série
d'éléments, où la position de chaque élément compte.\\

%\vspace{0.5 cm}
\begin{alertblock}{Syntaxe}

On peut accéder aux éléments d'une liste par leur \textbf{indice}.
\begin{verbatim}
tableau  : ["bus", "voiture", "vélo", "trotinette"]
indice :    0        1        2         3

\end{verbatim}
\end{alertblock}

\pause

\vspace{0.3 cm}
En Python, on utilise le type \texttt{\texttt{list}} pour les représenter. 
%En Python, une liste est créée en plaçant des éléments entre \textbf{crochets} \texttt{{[}{]}} et en les séparant par des \emph{virgules}.

\begin{exampleblock}{}
\small{
\begin{verbatim}
# Exemple: Liste des capitales européennes
capitales = ["Paris", "Berlin", "Madrid", "Athènes", "Zagreb", "Rome"]

# Exemple: Liste d'entiers
entiers = [5, 0, 17, 42, 23]
\end{verbatim}
}
\end{exampleblock}
\end{frame}

\frame[containsverbatim]{
\frametitle{Les listes}

 La construction de listes en Python est très \emph{flexible}.
\begin{itemize}
\item Il est possible de créer des listes contenant des valeurs de
\textbf{types différents}.
\end{itemize}

\begin{exampleblock}{}
\begin{verbatim}

# Liste mixte
mixte = [7, "tennis", True, 5.3]

\end{verbatim}
\end{exampleblock}
\small{Cette liste contient à la fois des valeurs de type \texttt{int},
\texttt{string}, \texttt{boolean} et \texttt{float}.}


\vspace{0.5 cm}


\begin{itemize}
\item La fonction \texttt{len()} nous permet de connaître le nombre
  d'éléments dans une liste.
\end{itemize}

\begin{exampleblock}{}
\begin{verbatim}
liste = [0, 1, 2, 3, 4, 5, 6, 7, 8, 9]
print("La liste contient", len(liste), "éléments.")
\end{verbatim}
\end{exampleblock}





}


\frame[containsverbatim]{
\frametitle{Longueur d'une liste}




\begin{itemize}
\item Une liste a une longueur variable. La méthode \texttt{append()} permet d'ajouter un élément à la fin d'une
liste.
\end{itemize}


\begin{exampleblock}{}
\begin{verbatim}
mixte = [7, "tennis", True, 5.3]

mixte.append("natation")
print(mixte)

\end{verbatim}
\end{exampleblock}

\begin{itemize}
\item  La méthode \texttt{remove()} permet d'enlever un élément  de la liste. 
\end{itemize}

\begin{exampleblock}{}
\begin{verbatim}

mixte.remove("tennis")
mixte.remove(7)
print(mixte)

\end{verbatim}
\end{exampleblock}

\small{Attention: les méthodes \texttt{append} et \texttt{remove} sont propres à l'objet (\texttt{mixte}). Chose à clarifier plus tard ...}

}




\frame[containsverbatim]{
\frametitle{Accéder aux éléments d'une liste}


 Pour une liste de longueur \(n\)   :


\begin{itemize}

\item   L'indice est la \textbf{position} d'un élément dans la liste.
\item  Pour une liste de \(n\) éléments, les indices vont de 0 à \(n-1\).
\item Les indices doivent être des nombres \textbf{entiers}.

 \end{itemize}

\begin{exampleblock}{Un premier exemple}
\begin{verbatim}
liste = ["bus", "voiture", "velo", "trotinette"]

print(liste[0])
print(liste[2])
\end{verbatim}
\end{exampleblock}


}

\frame[containsverbatim]{
\frametitle{Indiçage négatif}

 \small{Chaque élément d'une liste  a deux indices :
   un indice positif \(i\) et un indice négatif \(i - n\).}

\begin{alertblock}{}
\begin{verbatim}
liste          : ["bus", "voiture", "velo", "trotinette"]
indice positif :    0       1         2         3
indice négatif :   -4      -3        -2        -1 
\end{verbatim}
\end{alertblock}

\begin{exampleblock}{Indiçage négatif}
\begin{verbatim}
liste = ["bus", "voiture", "vélo", "trotinette"]
print(liste[-1])
print(liste[-2])
print(liste[0],liste[-4])
\end{verbatim}
\end{exampleblock}

\vspace{0.5 cm}

    \begin{itemize}

\item
  L'indiçage négatif revient à compter de la fin.
\item
  \textbf{Avantage} : Accèder au dernier ou avant-dernier élément de la
  liste, sans connaître sa longueur (indices -1 et -2 respectivement)
\end{itemize}

}


\frame[containsverbatim]{
\frametitle{Tranches}

On peut accéder à une partie (une tranche) d'une liste L de longueur $n$ en utilisant des indices :

\begin{itemize}
\item \texttt{L[m:k]} donne accès aux éléments \texttt{L[m], L[m+1], ... ,L[k-1]}
\item \texttt{L[m:]} donne accès aux éléments \texttt{L[m], ..., L[n-1]}
\item  \texttt{L[:k]} donne accès aux éléments \texttt{L[0], L[1], ... ,L[k-1]}  
\end{itemize}

\vspace{0.5 cm}
\begin{exampleblock}{Exemple}
\begin{verbatim}
liste = [2, 3, 5, 7, 11, 13, 17, 19, 21, 23]
print(liste[2:6])
print(liste[-8:-4])
liste[3:]
liste[:5]
liste[:-3]
\end{verbatim}
\end{exampleblock}

}


\frame[containsverbatim]{
\frametitle{Modifier ou ajouter des éléments dans une
liste}


\textbf{Rappel.} Les listes sont des objets \textcolor{blue}{mutables}. Leurs éléments peuvent
donc être modifiés.


\begin{itemize}
\item En écrivant \texttt{liste{[}i{]}\ =\ nouvel\_element}, l'élément de la
liste à l'indice \texttt{i} est replacé par \texttt{nouvel\_element}.
 \end{itemize}

\vspace{1 cm}

\begin{exampleblock}{}
\begin{verbatim}
capitales = ["Paris", "Berlin", "Madrid", "Athènes", "Zagreb"]
capitales[2:5] = ["Dublin", "Budapest", "Vienne"]
print(capitales)
\end{verbatim}
\end{exampleblock}





}




\frame[containsverbatim]{
\frametitle{Listes en compréhension}



\vspace{0.5 cm}

Imaginons, qu'on veut créer une liste contenant la racine carrée des
entiers \(0, 1, \dots, 9\).

\vspace{0.5 cm}

Avec une boucle: 
\begin{exampleblock}{}
\begin{verbatim}
import math
l = []
for i in range(10) :
    l.append(math.sqrt(i))

print(l)
\end{verbatim}
\end{exampleblock}

\vspace{0.5 cm}

De manière compacte, en une seule instruction:
\begin{exampleblock}{}
\begin{verbatim}
l = [math.sqrt(i) for i in range(10)]

print(l)

\end{verbatim}
\end{exampleblock}

}



\begin{frame}[fragile]{Listes en compréhension}

On peut créer des listes de façon \textbf{simple}, \textbf{concise} et
\textbf{compacte}.\\


\begin{alertblock}{Syntaxe}
%Une liste \texttt{L} en compréhension se créé de la façon suivante :
\texttt{L} = \texttt{{[}expr\ for\ x\ in\ t{]}}
\end{alertblock}

\begin{itemize}
\item
  Dans cette syntaxe : les crochets \texttt{{[}{]}}, les mots clés
  \texttt{for} et \texttt{in} sont obligatoires.
\item
  \texttt{t} est souvent une liste ou une collection de type
  \texttt{range}.
\item
  \texttt{x} est la variable de contrôle
\item
  \texttt{expr} est une expression qui dépend en général de \texttt{x}
  et dont la valeur est placée dans la liste.
\end{itemize}


\vspace{0.5 cm}

Exercice : Créer avec une seule instruction une liste qui contient la table de multiplication de 5.

\pause
\begin{exampleblock}{}
\begin{verbatim}
mult5 = [5*i for i in range(1,11)]
\end{verbatim}
\end{exampleblock}
\end{frame}

\begin{frame}[fragile]{Test d’appartenance}

On peut facilement tester l'appartenance d'un élément dans une liste avec l'opérateur \texttt{in}.

\begin{exampleblock}{}
\begin{verbatim}
nombres = [1, 3, 5, 7, 9, 11]
print(3 in nombres)
print(2 in nombres)
print(5 in nombres and 7 in nombres)
\end{verbatim}
\end{exampleblock}

\vspace{1 cm}

En combinant avec le mot-clé \texttt{not} on peut vérifier si un élément est absent d'une liste.

\begin{exampleblock}{}
\begin{verbatim}
4 not in nombres
\end{verbatim}
\end{exampleblock}

\end{frame}



\begin{frame}[fragile]{Copie d'une liste}

Qu'affiche-t-on à l'exécution? 

\begin{columns}
\column{4 cm}
\begin{exampleblock}{}
\begin{verbatim}
L1 = [1, 2, 3, 4]
L2 = L1
print(L2)

L2.append(5)
print(L2)
print(L1) 
\end{verbatim}
\end{exampleblock}
\column{5 cm}
\begin{tikzpicture}[
    font=\sffamily,
    frame/.style={draw, fill=gray!10, rounded corners, inner sep=1pt},
    object/.style={draw, fill=yellow!10, rounded corners, inner sep=0.5pt},
    arr/.style={->, thick, >=stealth, blue!60},
    node distance=0.3cm and 0.5cm
]

% Frames box
\node[frame, align=left, minimum width=0.3cm] (frame) at (3,0.5){
       \begin{tabular}{l }
        L1  \\
    \end{tabular}
};
\node[frame, align=left, minimum width=0.3cm] (frame1) at (3,1.5){
       \begin{tabular}{l }
        L2  \\
    \end{tabular}
};


\node[object, anchor=west] (Llist)  at (4,3.1) {
    \begin{tabular}{ccccc}        
        1 & 2 & 3 & 4 & \\
    \end{tabular}
};



% Arrows from L list entries to the row lists
%\draw[arr] ([yshift=1mm]Llist.north west) ++(0.3cm,-0.1) .. controls +(-0.5,1) and +(-0.5,-0.5) .. (list2.west);

% Arrow from L in frame to Llist
\draw[arr] ([xshift=0.6cm]frame.west)++(0.4cm,0) .. controls +(0.5,0.5) and +(-0.4,-0.4) .. (Llist.west);

\draw[arr] ([xshift=0.6cm]frame1.west)++(0.37cm,0) .. controls +(0.7,0.5) and +(-1,0) .. (Llist.west);

\only<2>
\node[object, anchor=west] (list3)  at (4,3.1) {
    \begin{tabular}{ccccccc}        
        1 & 2 & 3 & 4 & \textcolor{red}{5} \\
    \end{tabular}
};


\end{tikzpicture}



\end{columns}

}

\vspace{0.5 cm}

\begin{itemize}
\item
  Il ne faut jamais utiliser l'opérateur \texttt{=} pour copier une
  liste.
\item
  Les deux objets \texttt{L1} et \texttt{L2} partagent la \textbf{même
  zone mémoire}. Une modification de l'un entraine une modification de
  l'autre.
\end{itemize}

\end{frame}

\begin{frame}[fragile]{Copie d'une liste}


\begin{columns}
\column{4 cm}
\begin{exampleblock}{}
\begin{verbatim}
L1 = [1, 2, 3, 4]
L2 = L1[:]
print(L2)

L2.append(5)
print(L2)
print(L1)
\end{verbatim}
\end{exampleblock}
\column{5 cm}
\begin{tikzpicture}[
    font=\sffamily,
    frame/.style={draw, fill=gray!10, rounded corners, inner sep=1pt},
    object/.style={draw, fill=yellow!10, rounded corners, inner sep=0.5pt},
    arr/.style={->, thick, >=stealth, blue!60},
    node distance=0.3cm and 0.5cm
]

% Frames box
\node[frame, align=left, minimum width=0.3cm] (frame) at (3,0.5){
       \begin{tabular}{l }
        L2  \\
    \end{tabular}
};
\node[frame, align=left, minimum width=0.3cm] (frame1) at (3,1.5){
       \begin{tabular}{l }
        L1  \\
    \end{tabular}
};


\node[object, anchor=west] (Llist)  at (4.5, 1) {
    \begin{tabular}{cccc}        
        1 & 2 & 3 & 4  \\
    \end{tabular}
};


\node[object, anchor=west] (Llist1)  at (4.5,3.1.5) {
    \begin{tabular}{ccccc}        
        1 & 2 & 3 & 4 & \textcolor{red}{5}\\
    \end{tabular}
};


% Arrows from L list entries to the row lists
%\draw[arr] ([yshift=1mm]Llist.north west) ++(0.3cm,-0.1) .. controls +(-0.5,1) and +(-0.5,-0.5) .. (list2.west);

% Arrow from L in frame to Llist
\draw[arr] ([xshift=0.6cm]frame.west)++(0.4cm,0) .. controls +(0.5,0.5) and +(-0.4,-0.4) .. (Llist1.west);

\draw[arr] ([xshift=0.6cm]frame1.west)++(0.37cm,0) .. controls +(0.7,0.5) and +(-1,0) .. (Llist.west);


\end{tikzpicture}



\end{columns}
\end{frame}

\begin{frame}[fragile]{Copie d'une liste}

On peut utiliser la fonction $\texttt{list()}$ qui permet de créer une nouvelle liste. 


\begin{columns}
\column{5 cm}
\begin{exampleblock}{}
\begin{verbatim}
L1 = [1, 2, 3, 4]
L2 = list(L1)
print(L2)

L2.append(5)
print(L2)
print(L1)

\end{verbatim}
\end{exampleblock}
\column{5 cm}
\begin{tikzpicture}[
    font=\sffamily,
    frame/.style={draw, fill=gray!10, rounded corners, inner sep=1pt},
    object/.style={draw, fill=yellow!10, rounded corners, inner sep=0.5pt},
    arr/.style={->, thick, >=stealth, blue!60},
    node distance=0.3cm and 0.5cm
]

% Frames box
\node[frame, align=left, minimum width=0.3cm] (frame) at (3,0.5){
       \begin{tabular}{l }
        L2  \\
    \end{tabular}
};
\node[frame, align=left, minimum width=0.3cm] (frame1) at (3,1.5){
       \begin{tabular}{l }
        L1  \\
    \end{tabular}
};


\node[object, anchor=west] (Llist)  at (4.5, 1) {
    \begin{tabular}{cccc}        
        1 & 2 & 3 & 4  \\
    \end{tabular}
};


\node[object, anchor=west] (Llist1)  at (4.5,3.1.5) {
    \begin{tabular}{ccccc}        
        1 & 2 & 3 & 4 & \textcolor{red}{5}\\
    \end{tabular}
};


% Arrows from L list entries to the row lists
%\draw[arr] ([yshift=1mm]Llist.north west) ++(0.3cm,-0.1) .. controls +(-0.5,1) and +(-0.5,-0.5) .. (list2.west);

% Arrow from L in frame to Llist
\draw[arr] ([xshift=0.6cm]frame.west)++(0.4cm,0) .. controls +(0.5,0.5) and +(-0.4,-0.4) .. (Llist1.west);

\draw[arr] ([xshift=0.6cm]frame1.west)++(0.37cm,0) .. controls +(0.7,0.5) and +(-1,0) .. (Llist.west);


\end{tikzpicture}
\end{columns}

\end{frame}

\begin{frame}[fragile]{Concatener deux listes}

Qu'affiche-t-on à l'exécution? 

\vspace{0.5 cm}

\begin{exampleblock}{Exemple (rappel)}
\small{
\begin{verbatim}
capitales1 = ["Paris", "Berlin", "Madrid", "Athènes", "Zagreb", "Rome"]
capitales2 = ["Bruxelles", "Oslo"]

capitale1=capitales1 + capitales2
print(capitales1)  
\end{verbatim}
}
\end{exampleblock}

\vspace{0.5 cm}

Comment faire pour créer une nouvelle liste contenant la concatenation? 

\end{frame}


\begin{frame}[fragile]{Listes imbriquées}

Une liste où chaque élément est à nouveau une liste s'appelle 
\textbf{liste imbriquée}, \textbf{double liste} ou
\textbf{liste bidimensionnelle}.


\vspace{0.5 cm}

\begin{exampleblock}{}
\begin{verbatim}
L = [[1, 2, 3], [4, 5], [6, 7, 8, 9]]
print(L[0])
print(L[1])
print(L[2][1])
\end{verbatim}
\end{exampleblock}
\end{frame}


\begin{frame}[fragile]{Listes imbriquées}

On veut  calculer la somme de tous les éléments de la liste
\texttt{L}.


\begin{exampleblock}{}
\begin{verbatim}
L = [[1, 2, 3], [4, 5], [6, 7, 8, 9]]

somme = 0
for i in range(len(L)) :
    for j in range(len(L[i])) :
        somme += L[i][j]

print("La somme de tous les éléments de la liste vaut :", somme)
\end{verbatim}
\end{exampleblock}

\pause 

ou encore 


\begin{exampleblock}{}
\begin{verbatim}
somme = 0
for i in L :
    for j in i :
        somme += j
        
print("La somme de tous les éléments de la liste vaut :", somme)  
\end{verbatim}
\end{exampleblock}



\end{frame}

\begin{frame}[fragile]{Création de listes imbriquées}


\begin{exampleblock}{}
\begin{verbatim}
m = 3
n = 4
L = [[0] * n] * m
print(L)
\end{verbatim}
\end{exampleblock}
\texttt{m} listes sont créées mais qui font toutes référence à la même liste !

\vspace{0.5 cm}

\begin{columns}
\column{3 cm}
\begin{exampleblock}{}
\begin{verbatim}
L[0][1] = 1
print(L)
\end{verbatim}
\end{exampleblock}
%\vspace{0.5 cm}
%Comment faire de manière en évitant une copie de références ?
\column{5 cm}
\begin{tikzpicture}[
    font=\sffamily,
    frame/.style={draw, fill=gray!10, rounded corners, inner sep=1pt},
    object/.style={draw, fill=yellow!10, rounded corners, inner sep=0.5pt},
    arr/.style={->, thick, >=stealth, blue!60},
    node distance=0.3cm and 0.5cm
]

% Frames box
\node[frame, align=left, minimum width=0.3cm] (frame) at (3,0.5){
       \begin{tabular}{l }
        L  \\
    \end{tabular}
};
% List of lists (L)
\node[object, minimum width=0.5cm] (Llist)  at (4,2) {
    \begin{tabular}{c  |  c | c}
    $~ ^{\tiny{\texttt{0}}}$ & $~ ^{\tiny{\texttt{1}}}$ & $~ ^{\tiny{\texttt{2}}}$ \\     
    \end{tabular}
};

\node[object, anchor=west] (list2)  at (5,3.7) {
    \begin{tabular}{cccc}        
        0 & 0 & 0 & 0
    \end{tabular}
};


% Arrows from L list entries to the row lists
\draw[arr] ([yshift=1mm]Llist.north west) ++(0.3cm,-0.1) .. controls +(-0.5,1) and +(-0.5,-0.5) .. (list2.west);
\draw[arr] ([yshift=1mm]Llist.north west) ++(1cm,-0.1) .. controls +(-0.3,0.8) and +(-0.5,-0.5) .. (list2.west);
\draw[arr] ([yshift=1mm]Llist.north west) ++(1.5cm,-0.1) .. controls +(0,0.5) and +(-0.5,-0.5) .. (list2.west);

% Arrow from L in frame to Llist
\draw[arr] ([xshift=0.6cm]frame.west)++(0.25cm,0) .. controls +(0.5,0.5) and +(-1,0) .. (Llist.west);

\only<2>
\node[object, anchor=west] (list3)  at (5,3.7) {
    \begin{tabular}{cccc}        
        0 & \textcolor{red}{1} & 0 & 0
    \end{tabular}
};

\end{tikzpicture}
\end{columns}
\end{frame}

\begin{frame}[fragile]{Création de listes imbriquées}

Plusieurs façons correctes de faire. La plus compacte est sûrement la suivante :

\begin{exampleblock}{}
\begin{verbatim}
m = 3
n = 4
L = [[0] * n for i in range(m)]
print(L)
\end{verbatim}
\end{exampleblock}

\vspace{1 cm}

\begin{columns}
\column{6 cm}
\begin{tikzpicture}[
    font=\sffamily,
    frame/.style={draw, fill=gray!10, rounded corners, inner sep=1pt},
    object/.style={draw, fill=yellow!10, rounded corners, inner sep=0.5pt},
    arr/.style={->, thick, >=stealth, blue!60},
    node distance=0.3cm and 0.5cm
]

% Frames box
\node[frame, align=left, minimum width=0.3cm] (frame) at (3,0.5){
       \begin{tabular}{l }
        L  \\
    \end{tabular}
};
% List of lists (L)
\node[object, minimum width=0.5cm] (Llist)  at (4,2) {
    \begin{tabular}{c  |  c | c}
    $~ ^{\tiny{\texttt{0}}}$ & $~ ^{\tiny{\texttt{1}}}$ & $~ ^{\tiny{\texttt{2}}}$ \\     
    \end{tabular}
};

% Three 1D lists (rows)
\node[object, anchor=west] (list1) at (5,3) {
    \begin{tabular}{cccc}
             0 & 0 & 0 & 0
    \end{tabular}
};

\node[object, anchor=west] (list2)  at (5,3.7) {
    \begin{tabular}{cccc}        
        0 & 0 & 0 & 0
    \end{tabular}
};

\node[object, anchor=west] (list3)  at (5,4.4){
    \begin{tabular}{cccc}    
        0 & 0 & 0 & 0
    \end{tabular}
};

% Arrows from L list entries to the row lists
\draw[arr] ([yshift=1mm]Llist.north west) ++(0.3cm,-0.1) .. controls +(-0.5,1) and +(-0.5,-0.5) .. (list3.west);
\draw[arr] ([yshift=1mm]Llist.north west) ++(1cm,-0.1) .. controls +(-0.3,0.8) and +(-0.5,-0.5) .. (list2.west);
\draw[arr] ([yshift=1mm]Llist.north west) ++(1.5cm,-0.1) .. controls +(0,0.5) and +(-0.5,-0.5) .. (list1.west);

% Arrow from L in frame to Llist
\draw[arr] ([xshift=0.6cm]frame.west)++(0.25cm,0) .. controls +(0.5,0.5) and +(-1,0) .. (Llist.west);

\only<2->
\node[object, anchor=west] (list4) at (5,4.4) {
    \begin{tabular}{cccc}
             0 & \textcolor{red}{1} & 0 & 0
    \end{tabular}
};


\end{tikzpicture}


\column{3 cm}
\begin{exampleblock}{}
\begin{verbatim}
L[0][1] = 1
print(L)
\end{verbatim}
\end{exampleblock}
\end{columns}

\end{frame}






\begin{frame}[fragile]{Création de listes imbriquées}

\textbf{Exercice} : Créer une liste en deux dimensions, contenant la
matrice suivante :

\[\begin{bmatrix}1 & 2 & 2 & 2 \\0 & 1 & 2 & 2 \\ 0 & 0 & 1 & 2 \\ 0 & 0 & 0 & 1\end{bmatrix}\]

\pause

\begin{exampleblock}{}
\begin{verbatim}
L = [[0] * 4 for i in range(4)]
for i in range(4) :
    for j in range(4) :
        if i == j :
            L[i][j] = 1
        elif i < j :
            L[i][j] = 2
            
print(L)
\end{verbatim}
\end{exampleblock}

\end{frame}

\begin{frame}[fragile]{Création de listes imbriquées}

Solution plus compacte: 


\begin{exampleblock}{}
\begin{verbatim}
L = [0] * 4
for i in range(4) : # pour chaque ligne
    L[i] = [0] * i + [1] + [2] * (n - i - 1)

print(L)
\end{verbatim}
\end{exampleblock}

\end{frame}

\end{document}


\textbf{Remarque.} La fonction \texttt{len()} n'est pas propre aux listes
et peut s'appliquer à différents types d'objets.

\begin{exampleblock}{}
\begin{verbatim}
chaine = 'abracadabra'
print(len(chaine))
\end{verbatim}
\end{exampleblock}



\begin{frame}[fragile]{Chaînes de caractères}

\begin{exampleblock}{}
\begin{verbatim}
chaine = 'abracadabra'
chaine[2] = 'd'
print(chaine[2:5])
\end{verbatim}
\end{exampleblock}



\end{frame}



